% Constructed conceptual illustration, not an experiment or learned potential.
% Requires TikZ and arrows.meta; suitable for a full-width figure.
\begin{tikzpicture}[
  x=1cm,y=1cm,font=\sffamily\footnotesize,
  >={Latex[length=1.7mm,width=1.2mm]},
  frame/.style={draw=black!35,line width=.5pt,rounded corners=1.5pt},
  contour/.style={draw=black!17,line width=.55pt},
  field/.style={->,draw=black!40,line width=.6pt},
  update/.style={->,line width=1.05pt},
  smallnote/.style={font=\sffamily\scriptsize,text=black!70,align=center}
]
\node[anchor=west,font=\sffamily\bfseries\small] at (0,5.05)
  {(A) Time-conditioned field};
\node[anchor=west] at (0,4.60)
  {$U(x,t)$: the time input can change the update};
\node[anchor=west,font=\sffamily\bfseries\small] at (8.2,5.05)
  {(B) Time-unconditioned ReGAE};
\node[anchor=west] at (8.2,4.60)
  {$\Delta_z(x)=D_{\theta}(x,z)-x$: fixed $z$, no explicit $t$};
% Each thumbnail repeats the SAME state-space domain. They are snapshots
% at selected times, not a partition of states into different noise bands.
\foreach \shift/\time/\yc/\dy/\shade in
  {0/1/2.40/.85/blue,2.55/2/1.55/.02/teal,5.10/3/.75/-.70/violet}{
  \begin{scope}[shift={(\shift,1.05)}]
    \path[fill=\shade!2] (0,0) rectangle (2.15,2.92);
    \begin{scope}
      \clip (0,0) rectangle (2.15,2.92);
      \foreach \rx/\ry in {.32/.36,.67/.72,1.05/1.13,1.48/1.57}{
        \draw[draw=\shade!28,line width=.55pt]
          (1.60,\yc) ellipse[x radius=\rx cm,y radius=\ry cm];
      }
    \end{scope}
    \draw[frame] (0,0) rectangle (2.15,2.92);
    \node[font=\sffamily\small] at (1.075,3.20) {$t_{\time}$};
    \fill[black] (.54,1.40) circle[radius=1.5pt];
    \node[anchor=east,inner sep=2pt] at (.45,1.40) {$x$};
    \draw[update,draw=\shade!65!black]
      (.54,1.40)--(1.45,{1.40+\dy});
    % Identical coordinate indicators reinforce the common domain.
    \draw[field] (.17,.15)--(.63,.15);
    \draw[field] (.17,.15)--(.17,.59);
  \end{scope}
}
\node[smallnote] at (3.63,.60)
  {Same marked state $x$; time may be continuous.};
\draw[black!15,line width=.5pt] (7.75,.35)--(7.75,5.28);
% A single fixed code-conditioned state-update map. Contours are illustrative
% guides only; they do not imply that the map is an energy gradient.
\path[fill=blue!2] (8.2,1.05) rectangle (15.65,3.97);
\begin{scope}
  \clip (8.2,1.05) rectangle (15.65,3.97);
  \foreach \rx/\ry in {.40/.27,.83/.55,1.37/.87,1.99/1.15,2.66/1.48,3.35/1.82}{
    \draw[contour] (13.30,2.30) ellipse[x radius=\rx cm,y radius=\ry cm];
  }
  \draw[field] (10.05,3.36)--(10.76,3.18);
  \draw[field] (11.96,3.55)--(12.36,3.12);
  \draw[field] (14.91,1.55)--(14.42,1.86);
\end{scope}
\draw[frame] (8.2,1.05) rectangle (15.65,3.97);
% Two hand-constructed trajectories, with arrows marking repeated updates.
\draw[update,draw=blue!65!black] (8.95,1.48) to[out=4,in=205] (10.25,1.80);
\draw[update,draw=blue!65!black] (10.25,1.80) to[out=35,in=190] (11.40,2.50);
\draw[update,draw=blue!65!black] (11.40,2.50) to[out=10,in=165] (12.48,2.47);
\draw[update,draw=blue!65!black] (12.48,2.47)--(13.30,2.30);
\fill[blue!65!black] (8.95,1.48) circle[radius=1.8pt];
\node[smallnote,anchor=west,text=blue!65!black] at (8.66,1.22)
  {noise initialization};
\draw[update,draw=orange!80!black] (14.70,3.45) to[out=238,in=45] (14.20,2.94);
\draw[update,draw=orange!80!black] (14.20,2.94)--(13.75,2.59);
\draw[update,draw=orange!80!black] (13.75,2.59)--(13.30,2.30);
\fill[orange!80!black] (14.70,3.45) circle[radius=1.8pt];
\node[smallnote,anchor=east,text=orange!80!black] at (15.32,3.73)
  {perturbed state};
\filldraw[fill=black,draw=white,line width=.5pt]
  (13.30,2.30) circle[radius=2.3pt];
\node[smallnote,anchor=north] at (13.30,2.13)
  {illustrative target};
\node[smallnote] at (11.93,.60)
  {States evolve under one code-conditioned update map.};
\node[smallnote] at (7.83,-.02)
  {Illustrative fields; no learned energy surface is assumed.};
\end{tikzpicture}
